\section{Discussion}\label{sec:discussion}

Uncertainty in claims about repeated cross-national surveys is usually
computed for the wrong unit, twice. We have given a multiplicity-free
simultaneous band for the claim unit and a hard boundary for the calibration unit: the correction is
non-identified without the design-noise law (Theorem~\ref{thm:impossible}),
and unreachable for the deployed procedure at cross-national scale even with
it (Proposition~\ref{prop:unreach}). At the national unit the method reduces
to clustered population conformal, and the reported design share certifies
that the simple band is the right one. The deconvolution regime needs many
exchangeable units and comparable design noise, never both across countries
but reachable one unit down (Section~\ref{sec:evidence}).

\paragraph{Limitations.} The machinery prices sampling error only: measurement non-invariance
(wording changes, house effects, mode switches) produces the very shifts the
bands certify, so certifications straddling such changes are conditional on
measurement constancy. The
deployed $m$-of-$m$ PSU bootstrap \citep{rao1992recent,wolter2007introduction}
biases design variance downward; the Rao--Wu--Yue rescaled rerun leaves every
certification count and cross-national gate unchanged, while the small-area
activation set is bootstrap-sensitive in the disclosed direction
(Supplementary Material). Transport-band exchangeability is a
superpopulation postulate over a non-sampled, panel-attriting country set,
and holds worse for the post-communist bloc. And without design metadata
(the WVS trend file) only weights-only replication is possible.

\paragraph{The general lesson.} Attaching uncertainty to the right unit
changes substantive conclusions: no persistent decline certified among
thirty-three ESS countries over the full record, yet span erosion certified
in eight, mostly invisible to adjacent-wave readings; a certified
deconsolidation set several-fold smaller, concentrated outside the
consolidated West. The pattern is not confined to surveys: wherever prediction is calibrated on estimated objects
(Remark~\ref{rem:general}), the calibration targets carry their own error,
and the claim unit is often coarser than the trend asserted. The right
response is to state the strongest object the data support, correct
estimation noise where it is identified, and abstain where it is not, rather
than to widen every band reflexively.
